\section{Null spaces, column spaces, and linear transformations}
For an $m\times n$ matrix $A$, the null space $\Nul A$ contains the
\emph{inputs} $x\in\R^n$ satisfying $Ax=0$. The column space $\Col A$
contains the \emph{outputs} $Ax\in\R^m$. Do not confuse these two roles.
To test $w\in\Nul A$, multiply $Aw$ and compare with zero. To test
$w\in\Col A$, solve $Ax=w$ and ask whether a solution exists.
A homogeneous linear equation defines a subspace; a nonhomogeneous equation
usually describes a translated subspace and does not even contain zero.
For a linear map $T:V\to W$, the corresponding terms are kernel
$\{v:T(v)=0\}$ and range $\{T(v):v\in V\}$.

\practice{4.2.P1}{Two descriptions of one plane}
For $W=\{(a,b,c):a-3b-c=0\}$, first observe that
$W=\Nul[\,1\ -3\ -1\,]$. The null space theorem therefore makes it a
subspace of $\R^3$.
Alternatively solve $a=3b+c$. Then
\[
 (a,b,c)=b(3,1,0)+c(1,0,1).
\]
As $b,c$ vary freely, this proves
$W=\Span\{(3,1,0),(1,0,1)\}$, a subspace by the span theorem.
The first description is implicit (an equation); the second is explicit
(generators and arbitrary coefficients).

\practice{4.2.P2}{Add two attainable outputs}
Suppose $Ax=v$ and $Ay=w$ have solutions $x,y$. By distributivity,
$A(x+y)=Ax+Ay=v+w$, so the equation with right-hand side $v+w$ is
also consistent. Equivalently $v,w\in\Col A$, and the column space is
closed under addition. We need not find either of the original solutions.

\exercise{4.2.1}{Test a proposed null-space vector directly}
\goal Decide whether $w=(1,3,-4)$ is in $\Nul A$ for the displayed matrix.
Multiply using the original rows, retaining the signs:
\[
 \begin{pmatrix}3&-5&-3\\6&-2&0\\-8&4&1\end{pmatrix}
 \begin{pmatrix}1\\3\\-4\end{pmatrix}
 =\begin{pmatrix}3-15+12\\6-6+0\\-8+12-4\end{pmatrix}=0.
\]
Every coordinate is zero, so $w\in\Nul A$. A vector belongs to the null
space precisely when this multiplication gives zero, not merely when its
coordinates resemble a row or column of $A$.

\exercise{4.2.2}{Another direct multiplication}
With $w=(1,-1,1)$, compute
\[
 \begin{pmatrix}2&6&4\\-3&2&5\\-5&-4&1\end{pmatrix}
 \begin{pmatrix}1\\-1\\1\end{pmatrix}
 =\begin{pmatrix}2-6+4\\-3-2+5\\-5+4+1\end{pmatrix}=0.
\]
Thus $w\in\Nul A$. All three equations must hold; checking only the first
row would not have been enough.

\exercise{4.2.3}{Describe a null space using two free variables}
The homogeneous system for
$A=\left(\begin{smallmatrix}1&2&4&0\\0&1&3&-2\end{smallmatrix}\right)$ is
\[
 x_1+2x_2+4x_3=0,\qquad x_2+3x_3-2x_4=0.
\]
The pivot variables are $x_1,x_2$; choose $x_3=s,x_4=t$ freely.
The second equation gives $x_2=-3s+2t$, and substitution in the first gives
$x_1=-2(-3s+2t)-4s=2s-4t$. Therefore
\[
 x=s\begin{pmatrix}2\\-3\\1\\0\end{pmatrix}
   +t\begin{pmatrix}-4\\2\\0\\1\end{pmatrix}.
\]
Both parameters may be any real number, so the two displayed vectors span
$\Nul A$; this formula describes every solution and only solutions.

\exercise{4.2.4}{A forced zero coordinate and two free coordinates}
For $A=\left(\begin{smallmatrix}1&-3&2&0\\0&0&3&0\end{smallmatrix}\right)$,
the second equation is $3x_3=0$, so $x_3=0$. The first becomes $x_1=3x_2$.
Letting $x_2=s$ and $x_4=t$ gives
\[
 \Nul A=\left\{s\begin{pmatrix}3\\1\\0\\0\end{pmatrix}
 +t\begin{pmatrix}0\\0\\0\\1\end{pmatrix}:s,t\in\R\right\}.
\]
In particular the fourth column being zero leaves the fourth input free;
it does not require the fourth input to be zero.

\exercise{4.2.5}{Back substitution in a five-column matrix}
The nonzero rows give
\[
 x_1-4x_2+2x_4=0,\qquad x_3-5x_4=0,\qquad 2x_5=0.
\]
The last equation gives $x_5=0$. Set the free variables $x_2=s,x_4=t$.
Then $x_3=5t$ and $x_1=4s-2t$. Thus
\[
 \Nul A=\Span\left\{
 \begin{pmatrix}4\\1\\0\\0\\0\end{pmatrix},
 \begin{pmatrix}-2\\0\\5\\1\\0\end{pmatrix}\right\}.
\]
Each generating vector is obtained by giving one free variable value $1$
and the other value $0$; their arbitrary combinations preserve both free variables.

\exercise{4.2.6}{Three free variables}
The matrix gives two effective equations:
\[
 x_1+3x_2-4x_3-3x_4+x_5=0,\qquad x_2-3x_3+x_4=0.
\]
The zero third row adds no condition. Set $x_3=r,x_4=s,x_5=t$.
Then $x_2=3r-s$ and
$x_1=-3(3r-s)+4r+3s-t=-5r+6s-t$.
Consequently
\[
 \Nul A=\Span\left\{
 \begin{pmatrix}-5\\3\\1\\0\\0\end{pmatrix},
 \begin{pmatrix}6\\-1\\0\\1\\0\end{pmatrix},
 \begin{pmatrix}-1\\0\\0\\0\\1\end{pmatrix}\right\}.
\]

\exercise{4.2.7}{An affine plane misses zero}
The equation is $a+b+c=2$. At $(a,b,c)=(0,0,0)$ it would read $0=2$,
which is false. Hence the zero vector is not in $W$; $W$ is not a vector
space under the usual operations. The fact that the equation is linear in
the variables does not make its \emph{nonhomogeneous} solution set a subspace.

\exercise{4.2.8}{Move all variable terms to one side}
The relation $3r-2=3s+t$ becomes $3r-3s-t=2$.
Zero does not satisfy this relation, because its left side is zero.
Therefore the given set is not a vector space.

\exercise{4.2.9}{Two homogeneous constraints define a null space}
Rewrite $p-3q=4s$ and $2p=s+5r$ as
\[
 p-3q+0r-4s=0,\qquad 2p+0q-5r-s=0.
\]
Thus $W$ is the null space of
$\left(\begin{smallmatrix}1&-3&0&-4\\2&0&-5&-1\end{smallmatrix}\right)$,
and the null space theorem makes it a subspace of $\R^4$. Both right-hand
sides after rearrangement are zero; this is the essential property.

\exercise{4.2.10}{Another homogeneous description}
Move all terms to the left:
$3a+b-c+0d=0$ and $a+b+2c-2d=0$.
Then
\[
 W=\Nul\begin{pmatrix}3&1&-1&0\\1&1&2&-2\end{pmatrix}.
\]
The null space of a matrix is a subspace of its input space, here $\R^4$.
This proves the requested vector-space property without finding parameters.

\exercise{4.2.11}{A hidden constant obstructs the zero vector}
For $(s-2t,3+3s,3s+t,2s)$ to be zero, its fourth coordinate forces $s=0$.
But then its second coordinate is $3$, not zero. No choice of $t$ can fix
that contradiction. Thus zero is absent, and $W$ is not a vector space.
Merely setting all parameters zero would give only one candidate; the
argument above checks that \emph{no} parameters produce zero.

\exercise{4.2.12}{Another parametrization that never produces zero}
If $(3p-5q,4q,p,q+1)=0$, the third coordinate forces $p=0$ and the
second forces $q=0$. The fourth would then be $1$. This contradiction
shows that zero cannot occur, so $W$ is not a vector space.

\exercise{4.2.13}{A span in disguise}
For arbitrary real $c,d$,
\[
 (c-6d,d,c)=c(1,0,1)+d(-6,1,0).
\]
Both directions of this parametrization hold, so
$W=\Span\{(1,0,1),(-6,1,0)\}$. The span theorem makes it a subspace of $\R^3$.

\exercise{4.2.14}{A two-generator subspace}
Collecting the coefficients of $s,t$ gives
\[
 (-s+3t,s-2t,5s-t)=s(-1,1,5)+t(3,-2,-1).
\]
Since $s,t$ are unrestricted real numbers, $W$ is exactly the span of these
two vectors, hence is a subspace of $\R^3$.

\exercise{4.2.15}{Make the coefficient vectors into columns}
\goal Find a matrix whose column space is the given parametrized set.
Separate the coefficients of $r,s,t$:
\[
 \begin{pmatrix}2s+t\\r-s+2t\\3r+s\\2r-s-t\end{pmatrix}
 =r\begin{pmatrix}0\\1\\3\\2\end{pmatrix}
 +s\begin{pmatrix}2\\-1\\1\\-1\end{pmatrix}
 +t\begin{pmatrix}1\\2\\0\\-1\end{pmatrix}.
\]
Therefore one suitable matrix is
\[
 A=\begin{pmatrix}0&2&1\\1&-1&2\\3&1&0\\2&-1&-1\end{pmatrix}.
\]
By definition its column space contains exactly the linear combinations just
listed. Other matrices can have the same column space; uniqueness is not requested.

\exercise{4.2.16}{Construct a matrix from three parameters}
Here the required equality is
\[
 \begin{pmatrix}b-c\\2b+3d\\b+3c-3d\\c+d\end{pmatrix}
 =b\begin{pmatrix}1\\2\\1\\0\end{pmatrix}
 +c\begin{pmatrix}-1\\0\\3\\1\end{pmatrix}
 +d\begin{pmatrix}0\\3\\-3\\1\end{pmatrix}.
\]
Thus take
$A=\left(\begin{smallmatrix}1&-1&0\\2&0&3\\1&3&-3\\0&1&1\end{smallmatrix}\right)$.
For $x=(b,c,d)$, the product $Ax$ gives the original form, and every such
form is obtained from some $x$. This proves equality with $\Col A$.

\exercise{4.2.17}{Count inputs and outputs}
The displayed matrix has four rows and two columns, so its mapping is
$A:\R^2\to\R^4$.
\answer{(a) $\Nul A$ is a subspace of $\R^2$, so $k=2$.
(b) $\Col A$ is a subspace of $\R^4$, so $k=4$.}
These are ambient-space dimensions, not necessarily the dimensions of the two subspaces.

\exercise{4.2.18}{A rectangular matrix with three inputs}
The matrix has four rows and three columns. A vector $x$ in $Ax$ must
therefore have three coordinates, whereas the output has four.
\answer{(a) $k=3$ for the ambient space of $\Nul A$.
(b) $k=4$ for the ambient space of $\Col A$.}

\exercise{4.2.19}{Five inputs and two outputs}
The matrix has size $2\times5$, so $Ax$ is defined for $x\in\R^5$ and
belongs to $\R^2$.
\answer{(a) $\Nul A\subseteq\R^5$, so $k=5$.
(b) $\Col A\subseteq\R^2$, so $k=2$.}

\exercise{4.2.20}{A row matrix}
Here $A=[\,1\ -3\ 2\ 0\ -5\,]$ has one row and five columns.
It maps a five-coordinate input to one real number.
\answer{(a) $k=5$ for $\Nul A$. (b) $k=1$ for $\Col A$.}

\exercise{4.2.21}{Exhibit one vector in each space}
Using the matrix in Exercise 4.2.17, its second column equals $-2/3$ times
its first. Therefore $2a_1+3a_2=0$, giving the nonzero null vector $(2,3)$.
A direct check is
\[
 \begin{pmatrix}6&-4\\-3&2\\-9&6\\9&-6\end{pmatrix}
 \begin{pmatrix}2\\3\end{pmatrix}
 =\begin{pmatrix}12-12\\-6+6\\-18+18\\18-18\end{pmatrix}=0.
\]
The first column $(6,-3,-9,9)$ is a nonzero member of $\Col A$.
It equals $A(1,0)$, so its membership follows directly from the definition.

\exercise{4.2.22}{A null relation among three columns}
For the matrix in Exercise 4.2.18, choose $x=(-1,-1,1)$.
Multiplication verifies
\[
 Ax=\begin{pmatrix}-5+2+3\\1+0-1\\0+2-2\\5-7+2\end{pmatrix}=0.
\]
Thus this is a nonzero vector of $\Nul A$.
The first column $(5,-1,0,-5)$ is a nonzero vector in $\Col A$, because
it is the output $A(1,0,0)$.

\exercise{4.2.23}{A vector can belong to both spaces}
For $A=\left(\begin{smallmatrix}-2&4\\-1&2\end{smallmatrix}\right)$ and
$w=(2,1)$, the negative of the first column is $w$. Hence
$A(-1,0)=w$, proving $w\in\Col A$.
For null-space membership test the different multiplication:
\[
 Aw=(-4+4,-2+2)=(0,0).
\]
Thus $w\in\Nul A$ as well. These two statements are compatible; a square
matrix's input and output spaces have the same ambient space, though they
serve different roles.

\exercise{4.2.24}{Two independent membership checks in four dimensions}
The matrix and vector are
\[
 A=\begin{pmatrix}10&-8&-2&-2\\0&2&2&-2\\1&-1&6&0\\1&1&0&-2\end{pmatrix},
 \qquad w=\begin{pmatrix}2\\2\\0\\2\end{pmatrix}.
\]
\step{Column-space membership.}
The first two columns add to $w$, so $A(1,1,0,0)=w$.
This is a concrete solution of $Ax=w$, proving $w\in\Col A$.
\step{Null-space membership.}
Use $w$ itself as the input:
\[
 Aw=(20-16-4,\ 4-4,\ 2-2,\ 2+2-4)=(0,0,0,0).
\]
\answer{$w$ belongs to both $\Col A$ and $\Nul A$.}

\exercise{4.2.25}{Six statements about null spaces and column spaces}
\begin{enumerate}[label=(\alph*)]
\item \textbf{True.} $\Nul A=\{x:Ax=0\}$ is the definition.
\item \textbf{False.} An $m\times n$ matrix has $n$ input coordinates,
so its null space is in $\R^n$, not generally $\R^m$.
\item \textbf{True.} The possible outputs of $x\mapsto Ax$ are precisely
all combinations of the columns, hence form $\Col A$.
\item \textbf{False.} One consistent equation only says that this particular
$b$ is in the column space. For instance $A=\left(\begin{smallmatrix}1&0\\0&0\end{smallmatrix}\right)$
and $b=(1,0)$ give a consistent equation but $\Col A$ is just the horizontal axis.
\item \textbf{True.} The kernel contains zero and is closed under addition
and scalar multiplication because the map is linear.
\item \textbf{True.} $\Col A=\{Ax:x\in\R^n\}$; this is the output description
of the span of the columns.
\end{enumerate}

\exercise{4.2.26}{Six related statements about kernels and ranges}
\begin{enumerate}[label=(\alph*)]
\item \textbf{True.} A null space is a subspace of the input space, hence a vector space.
\item \textbf{True.} Columns of an $m\times n$ matrix have $m$ coordinates;
their combinations belong to $\R^m$.
\item \textbf{False.} Solutions of $Ax=b$ are inputs $x$, not outputs.
The column space instead contains those right-hand sides $b$ for which a solution exists.
\item \textbf{True.} The kernel of $x\mapsto Ax$ is, by definition, the set
of inputs taken to zero, exactly $\Nul A$.
\item \textbf{True.} The range is a subspace of the target vector space:
add two outputs or multiply an output by a scalar and linearity produces another output.
\item \textbf{True.} A homogeneous linear differential equation is
$L(f)=0$ for a linear differential operator $L$. Its solutions are the kernel
of $L$ on the vector space of functions with the required derivatives.
\end{enumerate}

\exercise{4.2.27}{Scale a known homogeneous solution}
The given solution vector is $u=(3,2,-1)$. The proposed vector is $10u$.
The solution set of the homogeneous system is a null space, which is a
vector space and is therefore closed under scalar multiplication.
Consequently $10u=(30,20,-10)$ is another solution. This uses no substitution
or row operations, as the exercise requests.

\exercise{4.2.28}{Scale an attainable right-hand side}
Both systems have the same coefficient matrix $A$. Their right-hand sides
are $b=(0,1,9)$ and $5b=(0,5,45)$, respectively.
Since the first is consistent, there is an input $u$ with $Au=b$.
Linearity gives $A(5u)=5Au=5b$, so the second system is consistent too.
Equivalently $b\in\Col A$ and column spaces are closed under scalar multiplication.

\exercise{4.2.29}{Prove the column-space subspace theorem}
A general member of $\Col A$ has form $Ax$, where $x\in\R^n$.
\step{(a) Zero is an output.}
Take $x=0$. Then $Ax=A0=0$, so the zero vector of $\R^m$ lies in $\Col A$.
\step{(b) Two outputs can be added.}
For inputs $x,w\in\R^n$, distributivity gives $Ax+Aw=A(x+w)$.
Since $x+w$ is another allowed input, this is an element of $\Col A$.
\step{(c) An output can be multiplied by any scalar.}
For real $c$, the identity $c(Ax)=A(cx)$ expresses it as an output as well.
\answer{The three conditions prove that $\Col A$ is a subspace of $\R^m$.}

\exercise{4.2.30}{The range of a linear transformation is a subspace}
Let $R=\{T(x):x\in V\}\subseteq W$. First $T(0)=0$, so $0\in R$.
For arbitrary $y=T(x)$ and $z=T(v)$ in $R$, linearity gives
\[
 y+z=T(x)+T(v)=T(x+v)\in R.
\]
For any real scalar $c$, $cy=cT(x)=T(cx)\in R$. The inputs $x+v,cx$
are in $V$ because $V$ is a vector space. Thus $R$ passes all subspace tests
inside $W$. One need not assume that $T$ is onto $W$.

\exercise{4.2.31}{Evaluate a quadratic at two inputs}
The map is $T(p)=(p(0),p(1))$ from $P_2$ into $\R^2$.
\step{(a) Check the two linearity identities at both coordinates.}
For polynomials $p,q$ and scalar $c$,
\[
 T(p+q)=(p(0)+q(0),p(1)+q(1))=T(p)+T(q),\qquad
 T(cp)=cT(p).
\]
Thus $T$ is linear.
\step{(b) Solve the kernel equations and then describe the range.}
Write $p(t)=a+bt+ct^2$. The equations $p(0)=p(1)=0$ give
$a=0$ and $b+c=0$, hence $p(t)=c(t^2-t)$.
The kernel is therefore $\Span\{t^2-t\}$.
For any desired output $(r,s)$, the polynomial $p(t)=r+(s-r)t$ satisfies
$p(0)=r$ and $p(1)=s$. Therefore the range is all of $\R^2$.
\answer{One kernel generator is $t^2-t$, and $T$ is onto $\R^2$.}

\exercise{4.2.32}{Repeat the same evaluation twice}
Here $T(p)=(p(0),p(0))$. With $p(t)=a+bt+ct^2$, the output is $(a,a)$.
It is zero exactly when $a=0$, leaving $b,c$ arbitrary. Thus
\[
 \ker T=\Span\{t,t^2\},\qquad
 \operatorname{range}T=\{(a,a):a\in\R\}=\Span\{(1,1)\}.
\]
The output has two coordinates, but those coordinates must be equal;
therefore the map is not onto the whole target $\R^2$.

\exercise{4.2.33}{Add a matrix to its transpose}
Define $T(A)=A+A^T$ on $M_{2\times2}$.
\step{(a) Verify linearity using transpose rules.}
For matrices $A,C$ and real scalar $r$,
\[
 T(A+C)=A+C+(A+C)^T=(A+A^T)+(C+C^T),\qquad
 T(rA)=r(A+A^T).
\]
These are exactly the two linearity identities.
\step{(b) Produce a preimage of any symmetric matrix.}
If $B^T=B$, choose $A=B/2$. Then
$T(A)=B/2+B^T/2=B$.
\step{(c) Establish equality of the range with the symmetric matrices.}
Every output satisfies $(A+A^T)^T=A^T+A=A+A^T$, so every output is symmetric.
Conversely part (b) realizes every symmetric matrix as an output. Both
inclusions prove the claimed equality of sets.
\step{(d) Solve the zero-output condition entry by entry.}
For $A=\left(\begin{smallmatrix}a&b\\c&d\end{smallmatrix}\right)$,
\[
 T(A)=\begin{pmatrix}2a&b+c\\b+c&2d\end{pmatrix}.
\]
This is zero exactly when $a=d=0$ and $c=-b$. Therefore
\[
 \ker T=\left\{b\begin{pmatrix}0&1\\-1&0\end{pmatrix}:b\in\R\right\}.
\]

\exercise{4.2.34}{The antiderivative with a fixed initial value}
For continuous $f$ on $[0,1]$, the specified antiderivative is uniquely
\[
 T(f)(x)=\int_0^x f(t)\,dt.
\]
Existence follows from the Fundamental Theorem of Calculus. For uniqueness,
two such antiderivatives differ by a constant, and their value zero at the
left endpoint forces that constant to be zero.
Linearity of the integral gives
$T(f+g)(x)=T(f)(x)+T(g)(x)$ and $T(cf)(x)=cT(f)(x)$ for every $x$.
Thus these are equalities of functions, and $T$ is a linear transformation.
If $T(f)$ is the zero function, differentiating it gives $f(x)=0$ on $(0,1)$,
and continuity gives zero at the endpoints too. Conversely the zero function
has zero antiderivative with the prescribed initial value.
\answer{The kernel consists only of the zero function.}

\exercise{4.2.35}{The image of a subspace}
Let $U$ be a subspace of $V$. Since $0\in U$ and $T(0)=0$, zero is in
$T(U)$. For $y=T(u)$ and $z=T(v)$ with $u,v\in U$, linearity gives
$y+z=T(u+v)\in T(U)$ because $u+v\in U$.
Likewise $cy=T(cu)\in T(U)$ since $cu\in U$ for every real $c$.
These three tests prove that $T(U)$ is a subspace of $W$.
The conclusion concerns the image of the whole subspace, not just a list of selected outputs.

\exercise{4.2.36}{The inverse image of a subspace}
Let $U=\{x\in V:T(x)\in Z\}$, where $Z$ is a subspace of $W$.
Since $T(0)=0\in Z$, zero belongs to $U$.
If $x,y\in U$, then $T(x),T(y)\in Z$, and
$T(x+y)=T(x)+T(y)\in Z$ by additive closure of $Z$. Thus $x+y\in U$.
For real $c$, $T(cx)=cT(x)\in Z$, so $cx\in U$.
Therefore $U$ is a subspace of $V$. This argument uses the preimage as a set;
it does not require $T$ to have an inverse function.

\exercise{4.2.37}{An output that is not a zero-output input}
The displayed data are
\[
 A=\begin{pmatrix}7&6&-4&1\\-5&-1&0&-2\\9&-11&7&-3\\19&-9&7&1\end{pmatrix},
 \qquad w=(1,1,-1,-3).
\]
To test null-space membership, multiply:
\[
 Aw=(7+6+4-3,\,-5-1+6,\,9-11-7+9,\,19-9-7-3)
 =(14,0,0,0)\ne0.
\]
Therefore $w\notin\Nul A$.
For the column space, the second column plus the third minus the fourth is
\[
 (6,-1,-11,-9)+(-4,0,7,7)-(1,-2,-3,1)=(1,1,-1,-3)=w.
\]
Thus $A(0,1,1,-1)=w$, which proves $w\in\Col A$.
\answer{$w$ is in the column space only.}

\exercise{4.2.38}{Another vector lying in both spaces}
Here
\[
 A=\begin{pmatrix}-8&5&-2&0\\-5&2&1&-2\\10&-8&6&-3\\3&-2&1&0\end{pmatrix},
 \qquad w=(1,2,1,0).
\]
First
$Aw=(-8+10-2,\,-5+4+1,\,10-16+6,\,3-4+1)=0$, so $w\in\Nul A$.
For the distinct column-space test use $x=(-2,-3,0,1)$:
\[
 Ax=(16-15,\,10-6-2,\,-20+24-3,\,-6+6)=(1,2,1,0)=w.
\]
\answer{$w$ belongs to both spaces, verified by two different products.}

\exercise{4.2.39}{Recover column relations and the entire null space}
The given matrix is
\[
 A=\begin{pmatrix}
 5&1&2&2&0\\3&3&2&-1&-12\\8&4&4&-5&12\\2&1&1&0&-2
 \end{pmatrix},\qquad B=[\,a_1\ a_2\ a_4\,].
\]
\step{(a) Exhibit the two column combinations.}
Coordinate-wise verification gives
\[
 a_3=\tfrac13a_1+\tfrac13a_2,\qquad
 a_5=\tfrac{10}3a_1-\tfrac{26}3a_2-4a_4.
\]
Both are combinations of columns of $B$, so $a_3,a_5\in\Col B$.
\step{(b) Solve all homogeneous equations.}
Exact row reduction gives
\[
 A\sim\begin{pmatrix}
 1&0&1/3&0&10/3\\0&1&1/3&0&-26/3\\0&0&0&1&-4\\0&0&0&0&0
 \end{pmatrix}.
\]
Take the free variables $x_3=s,x_5=t$. Then
$x_1=-s/3-10t/3$, $x_2=-s/3+26t/3$, and $x_4=4t$.
Hence
\[
 \Nul A=\Span\left\{
 \begin{pmatrix}-1/3\\-1/3\\1\\0\\0\end{pmatrix},
 \begin{pmatrix}-10/3\\26/3\\0\\4\\1\end{pmatrix}\right\}.
\]
These formulas also check the two column relations in part (a): moving
all terms to one side gives the corresponding zero combinations.
\step{(c) Explain both failures of invertibility separately.}
The first null-space generator is nonzero, so $T$ is not one-to-one: it
maps that generator and zero to the same output. There are only three
pivots among four rows, so the columns do not span $\R^4$; hence $T$ is
not onto $\R^4$. Its having five inputs does not by itself make it onto.

\exercise{4.2.40}{Find the line where two planes meet}
\goal Construct a nonzero vector in both spans and show that its span is the entire intersection.
Let $w=c_1v_1+c_2v_2=c_3v_3+c_4v_4$. Move the last two terms left and
solve the homogeneous system with columns $v_1,v_2,-v_3,-v_4$:
\[
 \begin{pmatrix}5&1&-2&0\\3&3&1&12\\8&4&-5&28\end{pmatrix}
 \sim
 \begin{pmatrix}1&0&0&-10/3\\0&1&0&26/3\\0&0&1&-4\end{pmatrix}.
\]
There is one free variable, $c_4=t$, and
$c_1=10t/3$, $c_2=-26t/3$, $c_3=4t$. Taking $t=1$ gives
\[
 w=\tfrac{10}3(5,3,8)-\tfrac{26}3(1,3,4)=(8,-16,-8).
\]
Independently the other plane gives
$4(2,-1,5)+(0,-12,-28)=(8,-16,-8)$, confirming both memberships.
Every possible set of coefficients in the homogeneous system is a multiple
of this one, so every common vector is a multiple of $w$.
\answer{$H\cap K=\Span\{(8,-16,-8)\}=\Span\{(1,-2,-1)\}$.}
